\documentclass[12pt,a4paper]{article}
\usepackage[utf8]{inputenc}
\usepackage[T1]{fontenc}
\usepackage[brazil]{babel}
\usepackage{geometry}
\usepackage{amsmath}
\usepackage{amssymb}
\usepackage{booktabs}
\usepackage{array}
\usepackage{xcolor}
\usepackage{listings}
\usepackage{hyperref}

\geometry{margin=2.5cm}

\hypersetup{
    colorlinks=true,
    linkcolor=blue!50!black,
    urlcolor=blue!60!black,
    citecolor=blue!50!black,
    pdftitle={Arquitetura Integrada de Deteccao de Queda com Central de Controle},
    pdfauthor={Equipe Inovacao UFLA}
}

\definecolor{codebg}{RGB}{248,248,248}
\definecolor{codeframe}{RGB}{220,220,220}

\lstdefinestyle{py}{
    language=Python,
    basicstyle=\ttfamily\small,
    keywordstyle=\color{blue!60!black},
    commentstyle=\color{gray!80!black},
    stringstyle=\color{teal!50!black},
    showstringspaces=false,
    breaklines=true,
    frame=single,
    rulecolor=\color{codeframe},
    backgroundcolor=\color{codebg},
    tabsize=4
}

\title{Arquitetura Integrada para Deteccao de Quedas\\\large Detection $\rightarrow$ Control Center $\rightarrow$ Notifier}
\author{Equipe Inovacao UFLA}
\date{14 de abril de 2026}

\begin{document}
\maketitle

\begin{abstract}
Este documento descreve a arquitetura integrada entre os projetos de reconhecimento de queda e notificacao, com uma camada intermediaria de decisao chamada \textbf{Central de Controle}. O objetivo principal e padronizar o fluxo de eventos e facilitar a extensao por novos metodos de entrada (novos sensores, novos algoritmos e novas fontes de evento), sem acoplar essas evolucoes ao modulo de notificacao.
\end{abstract}

\section{Visao Geral}
A arquitetura foi separada em tres responsabilidades:
\begin{enumerate}
    \item \textbf{Detection}: identifica eventos no dominio de entrada (queda, risco, etc.).
    \item \textbf{Control Center}: decide se um comando deve ser acionado, com politicas e cooldown.
    \item \textbf{Notifier}: executa o comando de notificacao para os canais externos.
\end{enumerate}

\subsection*{Fluxo principal}
\begin{center}
\fbox{\parbox{0.9\linewidth}{
\centering
\textbf{Detector de Entrada} $\Rightarrow$ \textbf{Evento Padronizado} $\Rightarrow$ \textbf{Central de Controle} $\Rightarrow$ \textbf{Comando de Notificacao} $\Rightarrow$ \textbf{Email/Telegram}
}}
\end{center}

\section{Estrutura dos Projetos}
\begin{itemize}
    \item \texttt{reconhecimento-queda/}: produtor de eventos de queda.
    \item \texttt{central-controle/}: modulo compartilhado de decisao e roteamento de comando.
    \item \texttt{notificador-quedas/}: envio de alertas para canais externos.
\end{itemize}

\section{Contrato de Evento}
Para permitir que outros colegas implementem novos metodos de input, a integracao usa um contrato simples de evento:

\begin{itemize}
    \item \texttt{event\_type}: tipo semantico do evento (ex.: \texttt{fall\_confirmed}).
    \item \texttt{source}: origem tecnica (ex.: \texttt{fall\_detection:skeleton}).
    \item \texttt{payload}: dados da ocorrencia (metadados e evidencias).
    \item \texttt{timestamp}: instante de geracao.
\end{itemize}

\subsection*{Exemplo de evento de queda}
\begin{lstlisting}[style=py]
ControlEvent(
    event_type="fall_confirmed",
    source="fall_detection:skeleton",
    payload={
        "subject": "Queda detectada",
        "body": "Foi detectada uma queda. Verifique o idoso.",
        "media_path": "capturas_quedas/queda_20260414-110101-123_4.jpg",
        "fall_id": 4,
        "detector": "skeleton",
        "confidence": 1.0,
        "min_confidence": 0.5,
    }
)
\end{lstlisting}

\section{Decisao na Central de Controle}
A Central de Controle recebe eventos e aplica politicas por tipo. No caso atual:

\begin{itemize}
    \item policy registrada para \texttt{fall\_confirmed};
    \item validacao de confianca minima;
    \item validacao de cooldown por comando (anti-disparo repetido).
\end{itemize}

Decisao simplificada:
\[
\text{disparar} = (\text{confidence} \ge \text{min\_confidence}) \land (\Delta t \ge \text{cooldown})
\]

\section{Saidas e Evidencias}
Quando uma queda e confirmada no detector:
\begin{enumerate}
    \item o detector registra no CSV;
    \item o detector salva snapshot em \texttt{capturas\_quedas/};
    \item o detector publica evento para a central;
    \item a central decide e, se aprovado, aciona o notificador.
\end{enumerate}

\section{Guia para Novos Metodos de Input}
Esta secao e o ponto principal para onboarding de colegas.

\subsection{Passo 1: Criar um produtor de evento}
Qualquer novo metodo (radar, IMU, audio, rede, regra clinica, modelo IA) deve produzir \texttt{ControlEvent}.

\subsection{Passo 2: Definir novo \texttt{event\_type}}
Exemplos:
\begin{itemize}
    \item \texttt{fall\_confirmed} (ja implementado)
    \item \texttt{prolonged\_inactivity}
    \item \texttt{abnormal\_gait}
    \item \texttt{manual\_panic\_button}
\end{itemize}

\subsection{Passo 3: Registrar policy para o novo tipo}
\begin{lstlisting}[style=py]
def prolonged_inactivity_policy(event: ControlEvent) -> ControlDecision:
    inactivity_seconds = float(event.payload.get("inactivity_seconds", 0))
    if inactivity_seconds < 120:
        return ControlDecision(False, "inatividade abaixo do limiar", command="notify_inactivity")

    return ControlDecision(True, "inatividade prolongada", command="notify_inactivity")

control_center.register_policy("prolonged_inactivity", prolonged_inactivity_policy)
\end{lstlisting}

\subsection{Passo 4: Reusar ou trocar gateway de saida}
Hoje e usado \texttt{SubprocessNotificationGateway} para chamar o notificador existente. No futuro, pode-se implementar outro gateway:
\begin{itemize}
    \item fila RabbitMQ/Kafka,
    \item API HTTP,
    \item log estruturado para SOC/NOC,
    \item integracao hospitalar (HL7/FHIR).
\end{itemize}

\section{Tabela de Responsabilidades}
\begin{center}
\renewcommand{\arraystretch}{1.25}
\begin{tabular}{>{\raggedright\arraybackslash}p{0.22\linewidth} >{\raggedright\arraybackslash}p{0.33\linewidth} >{\raggedright\arraybackslash}p{0.35\linewidth}}
\toprule
\textbf{Camada} & \textbf{Responsabilidade} & \textbf{Nao deve fazer} \\
\midrule
Detection & Detectar condicao de entrada e publicar evento & Enviar notificacao diretamente ao usuario final \\
Control Center & Aplicar regras de negocio, politica e cooldown & Conhecer detalhes internos de captura de camera \\
Notifier & Entregar mensagens para canais externos & Decidir regras clinicas ou logica de deteccao \\
\bottomrule
\end{tabular}
\end{center}

\section{Checklists para Implementacao}
\subsection*{Checklist de um novo input}
\begin{itemize}
    \item[ ] Definir novo \texttt{event\_type} e schema minimo de payload.
    \item[ ] Garantir envio de \texttt{source} identificavel e \texttt{timestamp}.
    \item[ ] Implementar policy e registrar na central.
    \item[ ] Testar casos: abaixo do limiar, acima do limiar, cooldown ativo.
    \item[ ] Validar comportamento com notificacao habilitada e desabilitada.
\end{itemize}

\subsection*{Checklist de qualidade para PR}
\begin{itemize}
    \item[ ] Nao quebrar contrato de evento existente.
    \item[ ] Evitar acoplamento do detector ao notificador.
    \item[ ] Mensagens de decisao claras para observabilidade.
    \item[ ] Documentar novo input neste documento.
\end{itemize}

\section{Execucao e Setup Recomendado}
No ambiente de desenvolvimento:
\begin{lstlisting}[style=py]
# na pasta central-controle
pip install -e .

# na pasta reconhecimento-queda
python src/fall_detection.py
\end{lstlisting}

No ambiente de producao, o recomendado e instalar \texttt{central-controle} como dependencia de versao fixa para evitar divergencias de API entre equipes.

\section{Conclusao}
A separacao em camadas cria uma base estavel para crescimento do sistema: novos metodos de input podem ser adicionados sem alterar o modulo de notificacao e sem reescrever a logica central de decisao. Com isso, o projeto fica mais colaborativo, escalavel e sustentavel para evolucao por multiplas equipes.

\end{document}
